\documentclass[10pt,a4paper]{amsart}
\usepackage[margin=19mm]{geometry}
\usepackage{amsmath,amssymb,mathtools}
\usepackage[scr=rsfs,cal=euler]{mathalfa}
\usepackage{xfrac}
\usepackage[unicode,hidelinks,bookmarks=false]{hyperref}
\newcommand{\ie}{i.e.\ }
\newcommand{\cf}{cf.\ }
\newcommand{\resp}{respectively\ }
\newcommand{\Subsection}{\subsection}
\providecommand{\nobreakdash}{\nobreak\hbox{-}}
\setlength{\emergencystretch}{2em}
\multlinegap0pt
\usepackage{bm}
\usepackage{bbm}
\usepackage{dsfont}
\usepackage{xspace}
\usepackage{cancel}
\usepackage{mathtools}
\usepackage{amsthm}
\makeatletter
\@ifundefined{theorem}{\newtheorem{theorem}{Theorem}}{}
\@ifundefined{lemma}{\newtheorem{lemma}[theorem]{Lemma}}{}
\makeatother
\newcommand{\dd}{\mathrm{d}}
\newcommand{\ii}{\mathrm{i}}
\title[Complex moments of the derivative of the Riemann zeta functi]{Complex moments of the derivative of the Riemann zeta function}
\author{Christopher Hughes, Andrew Pearce-Crump}
\date{Source: arXiv:2509.07788v1 (CC BY 4.0)}
\begin{document}
\maketitle

\noindent\emph{Source and licence.} Complex moments of the derivative of the Riemann zeta function. Version arXiv:2509.07788v1 (CC BY 4.0), distributed under the Creative Commons Attribution 4.0 International licence, as recorded in the arXiv licence metadata for this exact version and shown on the abstract page https://arxiv.org/abs/2509.07788v1. Full rights notice: licenses/arxiv-2509.07788-CC-BY-4.0.txt.\par\smallskip

\noindent\textit{Editorial scope.} Counted unit: the Lemma whose source label is \texttt{lem:FourierCoeffs} in the article (source0.tex lines 467--481, source label \texttt{lem:FourierCoeffs}), delivered together with the article's own complete original proof of it (source0.tex lines 492--544, one \texttt{proof} environment of 53 lines). No mathematics is written, repaired or re-derived: statement and proof are the article's own text line for line. The proof is detached from the statement in the source: the article states the result at lines 467--481 and prints its proof at lines 492--544, and the proof environment's own header names the statement, so the association is the article's own. Required context retained uncounted: U definition (lines 183--185); eq:Ftheta (lines 371--373). Every definition, display and label of those lines is retained unchanged. Omitted from this package and not redistributed: 1--182, 186--370, 374--466, 482--491, 545--758. Nothing inside a retained excerpt is omitted or altered: every retained body is delivered exactly as the article's lines are written, so no omission receipt of any kind accompanies this package. Citations: the retained excerpts carry no citation command at all, so no bibliography entry of the article is transcribed and no reference is redistributed. Reference adaptation: none was needed and none was made; every reference inside the delivered unit resolves inside it, so no unresolved reference or citation is produced. Printed slips kept literal: no printed word, symbol, hypothesis or number of the counted statement or of its proof was altered. Layout and class: the article's own class is \texttt{amsart} with packages including geometry, graphicx, placeins, natbib, xcolor, amsthm, amsmath, stmaryrd; the wrapper is the lane's pinned \texttt{amsart} editorial wrapper at 10pt on A4, which restates the theorem-like environments the retained text needs and adds the article's own macro definitions that the retained text needs (\texttt{\textbackslash dd}, \texttt{\textbackslash ii}), with extra wrapper packages (bm, bbm, dsfont, xspace, cancel, mathtools). The delivered numbering is the unit's own. Assets: no retained span contains a figure environment, a graphics-inclusion command, a PDF-page inclusion, a table or a TikZ picture; no image file is copied into this package, the package declares no asset dependency and no image file is redistributed with it. Licence: version arXiv:2509.07788v1 of this article is distributed under the Creative Commons Attribution 4.0 International licence, as recorded in the arXiv licence metadata for this exact version and shown on the abstract page https://arxiv.org/abs/2509.07788v1. A full rights notice is delivered at licenses/arxiv-2509.07788-CC-BY-4.0.txt. Characters: every retained excerpt is pure ASCII, so the delivered engine, XeLaTeX with its default Latin Modern text font, reports no missing character.
\begin{equation}\label{U definition}
U(z) = \int_1^e u(y) E_1(z\log y)\;\dd y \,,
\end{equation}

\begin{equation}\label{eq:Ftheta}
F_X(\vartheta) =  - \log\left(1-e^{-\ii \vartheta} \right) -\sum_{j=-\infty}^\infty U\big(\ii(\vartheta + 2\pi j)\log X\big)
\end{equation}

\begin{lemma} \label{lem:FourierCoeffs}
    Let $F_X(\vartheta)$ be defined in \eqref{eq:Ftheta}, then
    \[
    k F_X(-\vartheta) = \sum_{m=-\infty}^\infty s_m e^{\ii m \vartheta}
    \]
    where
    \[
    s_m =
    \begin{cases}
    0 & \text{ if } m \leq 0\\
    \frac{k}{m}  \left(1 - \int_1^{\exp(m/\log X)} u(y) \dd y \right) & \text{ if } 1 \leq m < \log X  \\
    0 & \text{ if } m\geq \log X
    \end{cases}
    \]
\end{lemma}

\begin{proof}[Proof of Lemma~\ref{lem:FourierCoeffs}]
Note that
\[
s_m = \frac{1}{2\pi} \int_{-\pi}^\pi k F_X(-\vartheta) e^{-\ii m \vartheta} \dd \vartheta .
\]
From \eqref{eq:Ftheta} we have
\[
k F_X(-\vartheta) =  - k\log\left(1-e^{\ii \vartheta} \right) -k \sum_{j=-\infty}^\infty U\big(\ii(-\vartheta + 2\pi j)\log X\big) .
\]
For $m$ an integer, we have
\begin{equation}\label{eq:FourierCoeffLog}
\frac{1}{2\pi} \int_{-\pi}^\pi -k \log\left(1-e^{\ii \vartheta} \right) e^{-\ii m \vartheta} \dd\vartheta =
\begin{cases}
    0 & \text{ if } m \leq 0\\
    \frac{k}{m} & \text{ if } m \geq 1
\end{cases}
\end{equation}
(To non-rigorously see why, note that $\log\left(1-e^{\ii \vartheta} \right) = -\sum_{\ell=1}^\infty e^{\ii \ell \vartheta}/\ell$ and the only term in the sum that survives the integration is when $\ell=m$ for positive integers $m$).

Furthermore, we have
\begin{multline*}
\frac{1}{2\pi}  \int_{-\pi}^\pi -k \sum_{j=-\infty}^\infty U\left(\ii(-\vartheta+2\pi j) \log X \right) e^{-\ii m \vartheta} \dd\vartheta
= \frac{-k}{2\pi}
\int_{-\infty}^\infty U(-\ii \vartheta \log X) e^{-\ii m \vartheta} \dd\vartheta\\
= \frac{-k}{2\pi} \int_{-\infty}^\infty \int_1^e u(y) E_1(-\ii \vartheta \log X\log y) e^{-\ii m \vartheta} \;\dd y \dd\vartheta
\end{multline*}
where the final equality comes from inserting the definition of $U$ from \eqref{U definition} and the support of $u$. Swapping the order of integration, this is
\[
\frac{-k}{2\pi}  \int_1^e u(y) \int_{-\infty}^\infty E_1(-\ii \vartheta \log X\log y) e^{-\ii m \vartheta} \; \dd\vartheta \dd y .
\]

The Fourier transform of the exponential integral is
\[
\frac{1}{2\pi} \int_{-\infty}^\infty E_1(-\ii \vartheta \log X\log y) e^{-\ii m \vartheta} \; \dd\vartheta  =
\begin{cases}
    0 & \text{ if } m<\log X \log y\\
    \frac{1}{2m} & \text{ if } m = \log X \log y\\
    \frac{1}{m} & \text{ if } m > \log X \log y
\end{cases}
\]

Therefore, making use of the fact that $u$ is supported on $[1,e]$ and has total mass 1, we have
\begin{multline*}
\frac{-k}{2\pi}  \int_1^e u(y) \int_{-\infty}^\infty E_1(-\ii \vartheta \log X\log y) e^{-\ii m \vartheta} \; \dd\vartheta \dd y = \frac{-k}{m} \int_1^{\max(1,\exp(m/\log X))} u(y) \dd y \\
=\begin{cases}
    0 & \text{ if } m \leq 0 \\
    -\frac{k}{m} \int_1^{\exp(m/\log X)} u(y) \dd y & \text{ if } 1 \leq m < \log X \\
    -\frac{k}{m}  & \text{ if } m \geq \log X
\end{cases}
\end{multline*}

The Fourier coefficient, $s_m$ is the sum of this and \eqref{eq:FourierCoeffLog}. Note that both terms are zero for $m\leq 0$ and the two terms perfectly cancel each other if $m\geq \log X$.
\end{proof}

\end{document}
